\section{Related Work}
\noindent \textbf{Diffusion Large Language Models.}
Autoregressive (AR) models have long been the dominant paradigm for natural language generation, largely due to the discrete and sequential nature of text. In contrast, diffusion models have achieved remarkable success in continuous domains such as image and video generation. Recently, diffusion-based language models have gained renewed interest and emerged as a competitive alternative for text generation, demonstrating promising progress across a wide range of tasks.

Representative approaches include pretraining dLLMs from scratch \cite{llada} and building dLLMs on top of existing AR models \cite{dream}. In parallel, several commercial systems \cite{seeddiffusion, mercury, gemini_diffusion} highlight the feasibility and practical potential of diffusion-based generation. Beyond these early successes, recent works \cite{fastdllmv2, llada2.0, wedlm} continue to advance dLLMs along multiple dimensions, including scaling to larger model size and exploring alternative training-inference paradigms for faster refinement. Moreover, diffusion-based modeling has been extended to multimodal settings, where multimodal dLLMs \cite{llada-v, mmada} demonstrate strong performance across a variety of tasks and point toward more unified diffusion-based generative models.

\noindent \textbf{Efficient Inference of dLLMs.}
Despite dLLMs's ability to update multiple positions in parallel, they still face practical challenges during inference, motivating more efforts on faster and more reliable decoding. 

Existing works focus on the KV Cache optimization \cite{fastdllm, dkvcache}, early stopping \cite{dllm-var, daedal}, distillation-based \cite{dparallel} acceleration, and others \cite{quantdllm, sparsedllm, dsb}. 
% Specifically, DSB \cite{dsb} uses a sliding block with a dynamic size to overcome the rigidity of the naive block, enabling efficient inference, which is orthogonal to \dawn{}. 
These directions have shown promising gains in improving the efficiency of dLLM inference. Beyond these optimizations, numerous studies have explored optimized sampling strategies for dLLM inference. Fast-dLLM \cite{fastdllm} adopts a confidence-aware strategy that unmasks multiple positions when their scores are sufficiently high, thereby making the independence approximation more reliable. EB-Sampler \cite{ebsampler} uses the entropy of predictive distributions to decide which positions are safe to update in parallel. KLASS \cite{klass} further incorporates temporal stability by comparing distributions across iterations via the KL divergence. 
WINO \cite{wino} follows a draft-and-verify style: it drafts many tokens in parallel and selectively regenerates those that fail verification. 
Spiffy and related works \cite{spiffy, orchestrating} apply speculative decoding style strategies to dLLMs to accelerate diffusion decoding. LocalLeap \cite{localleap} leverages a local determinism hypothesis, observing that positions adjacent to high-confidence commits tend to stabilize earlier and can therefore be updated more aggressively. 
% MEDAL \cite{medal} uses Monte Carlo tree search (MCTS) to look for better update trajectories among high-confidence candidates, and then fills in the remaining positions with a fast heuristic.
Unlike the above methods, this work explicitly approximates token coupling during inference and leverages this approximation to guide efficient parallel sampling.